% =========================================================
% Energy Conditions and Thermodynamic Consistency
% =========================================================

\section{Effective Stress-Energy Structure and Thermodynamic Consistency}

\subsection{Interpretive Status}

The representative TIG3 geometry introduced in the previous sections is not derived from a complete covariant modification of the Einstein field equations.

Accordingly, the effective stress-energy quantities introduced below should be interpreted conservatively as an effective geometric representation associated with the bounded-curvature vacuum sector.

The analysis is therefore interpretive and intended to probe the physical admissibility of the representative geometry rather than establish a complete matter-sector theory.

% =========================================================

\subsection{Effective Stress-Energy Tensor}

For the static spherically symmetric geometry

\begin{equation}
ds^2
=
-F(r)\,dt^2
+
\frac{dr^2}{F(r)}
+
r^2 d\Omega^2,
\end{equation}

with representative lapse function

\begin{equation}
F(r)
=
1
-
\frac{2Mr^2}{r^3+r_c^3},
\end{equation}

the geometric deviations from the Schwarzschild vacuum sector may be formally interpreted through an effective stress-energy tensor

\begin{equation}
T^{\mathrm{eff}}_{ab}.
\end{equation}

Within the static spherical sector, the effective structure may be parametrized by

\begin{equation}
T^{\mathrm{eff}\,a}_{\phantom{\mathrm{eff}}\,\,\,\,b}
=
\mathrm{diag}
\left(
-\rho_{\mathrm{eff}},
p_r,
p_t,
p_t
\right),
\end{equation}

where:
\begin{itemize}[leftmargin=0.6cm]
\item $\rho_{\mathrm{eff}}$ denotes the effective energy density,
\item $p_r$ the radial pressure,
\item and $p_t$ the tangential pressure.
\end{itemize}

The present interpretation is purely geometric and does not imply the existence of a fundamental physical fluid.

% =========================================================

\subsection{Effective Energy Density}

Using the Einstein tensor associated with the representative geometry, the effective energy density may be formally written as

\begin{equation}
\rho_{\mathrm{eff}}(r)
=
\frac{3Mr_c^3}
{4\pi (r^3+r_c^3)^2}.
\label{eq:rhoeff}
\end{equation}

Equation~(\ref{eq:rhoeff}) remains nonnegative for all

\begin{equation}
r \geq 0.
\end{equation}

Thus the representative exterior geometry satisfies the weak-energy-type positivity condition

\begin{equation}
\rho_{\mathrm{eff}} \geq 0.
\end{equation}

The present analysis should be interpreted conservatively as an effective admissibility check associated with the representative geometry rather than a complete energy-condition theorem.

% =========================================================

\subsection{Near-Core Effective Structure}

In the near-core regime

\begin{equation}
r \ll r_c,
\end{equation}

the representative lapse function behaves approximately as

\begin{equation}
F(r)
\approx
1
-
\frac{2M}{r_c^3}r^2.
\end{equation}

This structure resembles an effective de-Sitter-type core geometry with bounded curvature invariants.

Within this regime, the effective pressure sector may locally violate strong-energy-type conditions, similarly to previously studied regular black-hole geometries
\cite{bardeen1968,hayward2006,dymnikova1992,ansoldi2008}.

The present work interprets this conservatively as a bounded-curvature geometric effect associated with the representative vacuum structure.

No complete nonlinear stability theorem is established.

% =========================================================

\subsection{Surface Gravity and Hawking Temperature}

The surface gravity associated with the outer horizon radius $r_H$ is defined by

\begin{equation}
\kappa
=
\frac{1}{2}
F'(r_H).
\end{equation}

The corresponding Hawking temperature is

\begin{equation}
T_H
=
\frac{\kappa}{2\pi}
=
\frac{F'(r_H)}{4\pi}.
\label{eq:hawkingtemp}
\end{equation}

Using the representative TIG3 lapse function together with the horizon condition

\begin{equation}
F(r_H)=0,
\end{equation}

one obtains

\begin{equation}
T_H
=
\frac{1}{4\pi r_H}
\left(
1
-
\frac{3r_c^3}{r_H^3+r_c^3}
\right).
\label{eq:thfinal}
\end{equation}

% =========================================================

\subsection{Extremal Limit}

Near the critical branch-merger configuration associated with the TIG2 discriminant structure,

\begin{equation}
\beta \rightarrow \beta_c,
\end{equation}

the outer and inner horizon branches approach degeneracy.

In this regime,

\begin{equation}
F'(r_H)\rightarrow 0,
\end{equation}

and consequently

\begin{equation}
T_H \rightarrow 0.
\end{equation}

Thus the representative geometry admits a regular extremal-temperature limit without thermodynamic divergence within the analyzed sector.

The present work does not establish a complete black-hole thermodynamic framework beyond the semiclassical Hawking-temperature structure analyzed above.

% =========================================================

\subsection{Interpretation and Limitations}

The analysis developed in this section supports the interpretation that the representative TIG3 vacuum geometry possesses:

\begin{itemize}[leftmargin=0.6cm]
\item bounded-curvature representative structure,
\item nonnegative effective exterior energy density,
\item regular semiclassical temperature behavior,
\item and controlled extremal-limit structure.
\end{itemize}

However, several important questions remain open, including:

\begin{itemize}[leftmargin=0.6cm]
\item complete nonlinear stability,
\item full covariant field-equation closure,
\item exact matter-sector interpretation,
\item and nonperturbative thermodynamic consistency.
\end{itemize}

Accordingly, the present section should be interpreted as a conservative effective-consistency analysis associated with the representative TIG3 vacuum sector.
